\begin{annexitem}{annex:p19.bci.bayes}{thm:p19-bci-bayes}{Finite Bayes normalization and two different coin questions}
\begin{proof}
Put \(J(h,d)=\pi(h)L(d\mid h)\). Nonnegativity and
\(\sum_dL(d\mid h)=1\) imply
\(\sum_{h,d}J(h,d)=\sum_h\pi(h)=1\), so \(J\) is a joint
mass function. Its data marginal is \(m(d)\). For positive marginal,
conditioning gives \(J(h,d)/m(d)=q_d(h)\), and summing over \(h\)
gives one. If \(\pi(h)=0\), its numerator is zero. If \(m(d)=0\),
all joint masses at that datum are zero and the ratio supplies no conditional
law; any separate extension needs another convention rather than inference
by division.

For a fixed supported datum, let \(K(y\mid h,d)\) be nonnegative
and sum to one over the finite future alphabet for each \(h\). The
mixture is nonnegative and its sum is
\(\sum_hq_d(h)\sum_yK(y\mid h,d)=1\).

Each named five-toss string specifies five particular independent outcomes
of probability one half under the fair model. The product is \(2^{-5}=1/32\)
in both cases. Under the deterministic all-heads model only HHHHH has
positive mass, equal to one. With any positive prior weights on these
two models, observing HHTHT leaves only the fair model supported; observing
HHHHH does not. Their identical probability under the fair model is
therefore compatible with different posterior model odds. No biased-coin
distribution or participant interpretation is inferred from this construction.
\end{proof}
\end{annexitem}

\begin{annexitem}{annex:p19.bci.tree-count}{thm:p19-bci-tree-count}{Subset choices, internal clades and edge counts}
\begin{proof}
Each of the \(n\) objects is independently either included in or excluded
from a subset. This gives a bijection between subsets and \(n\)-bit
membership strings, hence \(2^n\) subsets.

Let \(I\) be the number of internal nodes in the rooted full binary
tree. Every internal node has two children, giving \(2I\) edges.
A finite tree with \(I+n\) nodes has \(I+n-1\) edges: starting
at the root, each other node has exactly one incoming parent edge.
Thus \(2I=I+n-1\), so \(I=n-1\), and the edge count is
\(2n-2\). Each internal node has at least two descendant leaves.

Internal nodes give distinct descendant-leaf sets. If one is a strict
ancestor of the other, its other child contributes at least one further
leaf; if neither is an ancestor, their descendant sets are disjoint.
Thus restricting concepts to internal clades gives exactly \(n-1\)
nonsingleton concepts. Leaf clades add \(n\) different singletons,
giving \(2n-1\) clades in total. The empty set is not one of them.
These counts use the stated full-binary and candidate conventions, not
an ambiguous claim about all possible branches or all real categories.
\end{proof}
\end{annexitem}

\begin{annexitem}{annex:p19.bci.iterated}{thm:p19-bci-iterated}{Stationarity, geometric convergence and failure cases}
\begin{proof}
Strict positivity gives \(m(d)>0\) for all data and positive posterior
entries. For each current hypothesis \(j\),
\[
 \sum_iQ_{ij}=\sum_dL(d\mid j)\sum_iq_d(i)=1.
\]
Thus \(Q\) is column stochastic and all its entries are positive.
The posterior formula yields
\[
 \pi(j)Q_{ij}
 =\pi(j)\pi(i)\sum_d\frac{L(d\mid i)L(d\mid j)}{m(d)}
 =\pi(i)Q_{ji}.
\]
Summing over \(j\) gives
\((Q\pi)_i=\pi(i)\sum_jQ_{ji}=\pi(i)\), because column
\(i\) sums to one. This establishes stationarity, not yet attraction.

For a direct convergence proof set
\(a_i=\min_jQ_{ij}>0\) and
\(\varepsilon=\sum_i a_i\). Any column sum bounds that total
by one, so \(0<\varepsilon\leq1\). Put \(r_i=a_i/\varepsilon\).
If \(\varepsilon=1\), every column must equal \(r\); thus one
step sends every input distribution to \(r=\pi\).
Otherwise define
\[
 R_{ij}=\frac{Q_{ij}-\varepsilon r_i}{1-\varepsilon}.
\]
Its entries are nonnegative and its columns sum to one. For any real vector
\(v\), column stochasticity and the triangle inequality give
\[
 \lVert Rv\rVert_1
 \leq\sum_{i,j}R_{ij}|v_j|=\sum_j|v_j|=\lVert v\rVert_1.
\]
For two probability columns \(\mu,\nu\), their difference sums
to zero. The shared \(\varepsilon r\) part cancels, hence
\[
 \lVert Q\mu-Q\nu\rVert_1
 \leq(1-\varepsilon)\lVert\mu-\nu\rVert_1.
\]
Using stationary \(\nu=\pi\) and iterating proves, for positive integers \(k\),
\[
 \lVert Q^k\mu-\pi\rVert_1
 \leq(1-\varepsilon)^k\lVert\mu-\pi\rVert_1\longrightarrow0.
\]
It also proves uniqueness of the stationary distribution: two stationary
probability columns would contradict strict contraction unless equal.

The produced-data masses after hypothesis law \(\mu\) are
\(s_\mu(d)=\sum_hL(d\mid h)\mu(h)\). Applying the same
nonnegative-sum inequality yields
\(\sum_d|s_\mu(d)-m(d)|\leq\lVert\mu-\pi\rVert_1\).
Consequently they converge to the prior predictive law when hypothesis
laws converge to \(\pi\). Alternately drawing data and hypotheses
uses the two conditional laws of \(J(h,d)=\pi(h)L(d\mid h)\),
which is the declared Gibbs correspondence; no human algorithm is thereby
identified.

For the MAP failure, take hypotheses \(A,B\), prior masses .9 and .1,
and a single datum of likelihood one under either hypothesis. The posterior
always equals the prior, but its unique mode is \(A\). MAP transmission
therefore produces point mass on \(A\), not the prior mixture.

For the nonmixing failure allow zeros: let \(A\) generate only datum
\(a\) and \(B\) only \(b\), with both prior masses positive.
The posterior identifies the generating hypothesis exactly, so \(Q\)
is the identity matrix. The prior is stationary, but every other initial
distribution is also stationary and need not converge to it. Strict positivity
is sufficient, not asserted necessary; more general ergodic conditions
can also give convergence. The two failures demonstrate why the source's
summary cannot be interpreted as convergence from Bayes's identity alone.
\end{proof}
\end{annexitem}
